\begin{abstract}
Robotic pollination in greenhouse settings needs flower detectors that adapt to a new crop with very few labels. We study a domain-matched alternative to standard COCO transfer: pretrain YOLOv8n jointly on three source crops (apple, strawberry, tomato) with labels harmonized to a single \texttt{flower} class, then fine-tune on a held-out target crop (kiwi) at five label budgets between 10 and 181. Naive fine-tuning of the joint pretrain underperforms COCO transfer at most budgets and collapses entirely at 25 labels (mAP@0.5 of 0.39 versus COCO's 0.89), driven by small-batch gradient noise that destroys the pretrain's flower-specific features. Freezing the YOLOv8n backbone during fine-tuning eliminates the collapse and produces the strongest results we have at the smallest budgets: at $n{=}10$ labels CrossPoll with backbone freezing reaches 0.630 mAP@0.5 versus 0.362 for COCO transfer, with roughly five times lower variance across seeds; at $n{=}25$ the two methods are within two points of each other. From $n{=}50$ onward COCO transfer regains a small lead. The result is a budget-dependent recommendation: for the realistic greenhouse-pollination scenario where an operator can label only a handful of images for a new crop, joint multi-crop pretraining with backbone freezing is the better protocol. We release the code, the harmonized dataset specification, the joint pretrain checkpoint, and a single \texttt{summary.csv} from which the table and figure in this paper are generated.
\end{abstract}
